\documentclass[10pt,a4paper]{amsart}
\usepackage[margin=19mm]{geometry}
\usepackage{amsmath,amssymb,mathtools}
\usepackage[scr=rsfs,cal=euler]{mathalfa}
\usepackage{xfrac}
\usepackage[unicode,hidelinks,bookmarks=false]{hyperref}
\newcommand{\ie}{i.e.\ }
\newcommand{\cf}{cf.\ }
\newcommand{\resp}{respectively\ }
\newcommand{\Subsection}{\subsection}
\providecommand{\nobreakdash}{\nobreak\hbox{-}}
\setlength{\emergencystretch}{2em}
\multlinegap0pt
\usepackage{amssymb}
\newtheorem{proposition}{Proposition}
\newtheorem{theorem}{Theorem}
\def\B{2}
\def\Om{\Omega}
\def\al{\alpha}
\def\bbset#1:#2\eeset{\{#1\,:\,#2\}}
\def\be{\beta}
\def\cJ{{\mathcal J}}
\def\cM{{\mathcal M}}
\def\cl#1{\overline{#1}}
\def\diag{\mathop{\bigtriangleup}}
\def\es{\varnothing}
\def\fun{{}^\frown}
\def\ga{\gamma}
\def\ilcc[#1,#2]{\ilsymlc#1\ilcomma/ #2\ilsymrc}
\def\ilcomma/{\boldsymbol{,}}
\def\ilcp[#1,#2){\ilsymlc#1\ilcomma/ #2\ilsymrp}
\def\ilpc(#1,#2]{\ilsymlp#1\ilcomma/ #2\ilsymrc}
\def\ilpp(#1,#2){\ilsymlp#1\ilcomma/ #2\ilsymrp}
\def\ilsymlc{\boldsymbol{[}}
\def\ilsymlp{\boldsymbol{(}}
\def\ilsymrc{\boldsymbol{]}}
\def\ilsymrp{\boldsymbol{)}}
\def\jumps#1{\mathfrak{J}(#1)}
\def\ka{\kappa}
\def\la{\lambda}
\def\lex#1#2{\mathfrak{L}(#2^{#1})}
\def\lexab{\lex\la\B}
\def\rarr{\Rightarrow}
\def\set#1{\bbset#1\eeset}
\def\sq#1#2{(#1)_{#2}}
\def\sset#1{\{#1\}}
\def\te{\theta}
\title{Homogeneous linearly ordered spaces}
\author{Anton Lipin, Evgenii Reznichenko}
\date{Source: arXiv:2609.15499v1 (CC BY 4.0)}
\begin{document}
\maketitle

\noindent\emph{Source and licence.} Homogeneous linearly ordered spaces. Version arXiv:2609.15499v1 (CC BY 4.0), distributed by arXiv under the Creative Commons Attribution 4.0 International licence, http://creativecommons.org/licenses/by/4.0/, as recorded in the arXiv licence metadata for this exact version and shown on the abstract page https://arxiv.org/abs/2609.15499v1. Full licence notice: \texttt{licenses/arxiv-2609.15499-CC-BY-4.0.txt}.\par\smallskip

\noindent\textit{Editorial scope.} Counted unit: the article's theorem t:blex:t1 (source0.tex lines 1370--1385). The article's complete original proof of it is delivered with it (source0.tex lines 1386--1480, one \texttt{proof} environment of 95 lines). No mathematics is written, repaired or re-derived: the statement and the proof are the article's own lines, in the article's own notation, and no retained line is substituted or re-flowed.  Retained uncounted as the source-local context the proof needs: lines 1347--1350; lines 1364--1366.  Nothing was omitted from any retained span, so the package carries no omission record.  No retained line contains a citation, so the wrapper carries no bibliography and no bibliography entry.  No retained line contains a figure environment, an image-inclusion command or drawing code, so no image is delivered and the package declares no asset dependency.  Editorial formatting only, in addition: an AMS \texttt{amsart} preamble that restates the article's own declarations and macros used by the retained lines, namely the environments \texttt{proposition}, \texttt{theorem} on separate counters, together with the standard packages those lines need.  The retained environments print in this document as Proposition 1, Theorem 1 in order of appearance, whereas the article prints the same items with the numbering of its own full document.  The complete native source of the article as distributed in the arXiv e-print for this exact version is bundled unchanged as \texttt{sources/210422/source0.tex}.
\begin{proposition}\label{p:blex:2-1}
Let $f: X\to Y$ be a surjective increasing map between LOTS $X$ and $Y$.
If $f^{-1}(y)$ is a closed interval in $X$ for every $y\in Y$, then $f$ is continuous.
\end{proposition}

\begin{proposition}\label{p:blex:2}
Let $\ka\leq \la$ be ordinals. The projection $\pi^\la_\ka: \lexab \to \B^\ka$ is continuous.
\end{proposition}

\begin{theorem}\label{t:blex:t1}
Let $X$ be a zero-dimensional compact LOTS, $\la$ a regular cardinal.
The following conditions are equivalent:
\begin{enumerate}
\item
$X$ and $\lexab$ are order isomorphic;
\item
$X$ and $\lexab$ are homeomorphic;
%\item
%$\chi(x,X)=\la$ for all $x\in X$ and $X$ embeds in a product $\prod_{\al<\la}X_\al$ of spaces such that $\chi(X_\al)<\la$ for $\al<\la$;
\item
$\chi(x,X)=\la$ for all $x\in X$ and $X$ embeds in a product $\prod_{\al<\la}X_\al$ of spaces such that $\chi(X_\al)<\la$ for $\al<\la$;
\item
$\chi(x,X)=\la$ for all $x\in X$ and the set $M$ of jump points of the LOTS $X$ decomposes into a union $M=\bigcup_{\al<\la} M_\al$ of subsets such that $\chi(\cl{M_\al})<\la$ for $\al<\la$.
\end{enumerate}
\end{theorem}

\begin{proof}
(1) $\rarr$ (2). Obvious.

(2) $\rarr$ (3). By Proposition \ref{p:blex:2}, $\lexab$ embeds via the map $\diag_{\al<\la} \pi_\al^\la$ into the product $\prod_{\al<\la} \lex\al\B$. It remains to note that $\chi(\lex\al\B)\leq |\al|<\la$ for $\al<\la$.

(3) $\rarr$ (4). Assume $X\subset \prod_{\al<\la}X_\al$. We may assume that each $X_\al$ is compact. Denote $P_\be = \prod_{\al<\be}X_\al$ for $\be\leq\la$ and let $p_\be$ be the projection from $P_\la$ onto $P_\be$ for $\be<\la$.

Let $\be<\la$. Since $\la$ is a regular cardinal and $\chi(X_\al)<\la$ for $\la<\be$, we have $\ka'=\sup\{ \chi(X_\al) : \al < \be \}<\la$. Set $\ka=\ka' + |\be|$. Then $\chi(P_\be)\leq\ka<\la$.

For $(a,b)\in\jumps X$ set
\[
\te_{a,b}=\min \set{\al<\la: p_\al(\ilpc(-\infty,a]) \cap  p_\al(\ilcp[b,+\infty))=\es}.
\]
Set $M_\be=\set{a,b : (a,b)\in\jumps X,\ \te_{a,b}=\be}$ for $\be<\la$. Then $M=\bigcup_{\be<\la} M_\be$. Let us show that $\chi(\cl{M_\be})<\la$. Suppose otherwise. Then $\chi(x,\cl{M_\be})=\la$ for some $x\in \cl{M_\be}$. Let $L=\ilpp(-\infty,x)\cap M_\be$ and $R=\ilpp(x,+\infty)\cap M_\be$. Then $\chi(x,\cl{L})=\la$ or $\chi(x,\cl{R})=\la$. Assume, without loss of generality, that $\chi(x,\cl{R})=\la$. Let $\ka=\chi(p_\be(x), P_\be)$ and $F=p_\be^{-1}(p_\be(x))$. Since $\chi(x,\cl{R})=\la$, $F$ is a set of type $G_\ka$ and $\ka<\la$, there is $y>x$ with $\ilpp(x,y)\subset F$, and the set $\ilpp(x,y)\cap M_\be$ is infinite. There is $(a,b)\in\jumps X$ with $\sset{a,b}\subset \ilpp(x,y)\cap M_\be$. Then $\sset{a,b}\subset F$, and hence $p_\be(a)=p_\be(b)=x$, contradicting $p_\be(\ilpc(-\infty,a]) \cap  p_\be(\ilcp[b,+\infty))=\es$.

(4) $\rarr$ (1). Set $\cM_\al=\set{(a,b)\in \jumps X: \sset{a,b}\cap M_\al\neq\es}$ for $\al<\la$ and define the map
\[
\Om: \jumps X \to \la,\ (a,b) \mapsto \min\set{\al<\la: (a,b)\in \cM_\al}.
\]
Set $\te_{u,v}=\min  \Om(\jumps{\ilcc[u,v]})$ and $\cJ_{u,v}=\Om^{-1}(\te_{u,v})\cap \jumps{\ilcc[u,v]}$ for $u,v\in X$, $u<v$. Fix $(a_{u,v},b_{u,v})\in \cJ_{u,v}$ for every $u,v\in X$, $u<v$.

By induction on $\al<\la$, we construct increasing continuous maps $f_\al: X\to \B^\al$ such that the following conditions hold:
\begin{enumerate}
\item
$f_\al=\pi_\al^\be\circ f_\be$, where $\al<\be<\la$ and $\pi_\al^\be: \B^\be\to \B^\al$ is the projection;
%\item if $\al<\la$ is a limit ordinal, then 
\item
if $\al<\la$, 
$s\in \B^\al$ and $\ilcc[u,v]=f^{-1}_\al(s)$, then $f^{-1}_{\al+1}(s\fun 0)=\ilcc[u,a_{u,v}]$ and $f^{-1}_{\al+1}(s\fun 1)=\ilcc[b_{u,v},v]$.
\end{enumerate}
Base of the induction, $\al=0$. Then $\B^\al=\sset\es$; let $f_\al: X\to \sset\es$ be the identity map.

Inductive step, $0<\al<\la$. Consider the case when $\al$ is a limit ordinal. For $x\in X$ define $f_\al(x)\in \B^\al$: $f_\al(x)(\be)=f_{\be+1}(x)(\be)$ for $\be<\al$. Consider the case when $\al$ is a non-limit ordinal, $\al=\be+1$.
Define the map $g_\al: X\to \B$. Let $x\in X$. Since $\chi(x,X)=\la$ and $\chi(\B^\al)<\la$, we have $u<v$, where $\ilcc[u,v]=f^{-1}_\be(f_\be(x))$. Set
\[
g_\al(x) = \begin{cases}
0,&	x\in \ilcc[u,a_{u,v}],
\\
1,&	x\in \ilcc[b_{u,v},v].
\end{cases}
\]
Define $f_\al(x)\in \B^\al$:
\[
f_\al(x)(\ka) = \begin{cases}
f_\be(x)(\ka),&	\ka<\be,
\\
g_\al(x),&	\ka=\be.
\end{cases}
\]
The maps $f_\al$ for $\al<\la$ have been constructed.

Define the map $f: X\to \lexab$. For $x\in X$ set $f(x)(\al)=f_{\al+1}(x)(\al)$.
Since each $f_\al$ is a monotone map, $f$ is a monotone map.

For $s\in \lexab$ and $\al<\la$ denote 
\begin{align*}
\ilcc[u_\al^s,v_\al^s]&=f^{-1}_\al(\pi_\al^\la(s)), 
&
(a_\al^s,b_\al^s)&=(a_{u_\al^s,v_\al^s},b_{u_\al^s,v_\al^s}).
\end{align*}
Then
\begin{align*}
\ilcc[u_{\be+1}^s,v_{\be+1}^s] &=
\begin{cases}
\ilcc[u_{\be}^s,a_{\be}^s], &  s(\be)=0,
\\
\ilcc[b_{\be}^s,v_{\be}^s], &  s(\be)=1,
\end{cases}
\\
\ilcc[u_{\ga}^s,v_{\ga}^s] &= \bigcap_{\al<\ga}\ilcc[u_\al^s,v_\al^s],
\end{align*}
for $\be<\la$ and limit $\ga<\la$. Hence 
\[
f^{-1}(s) = \ilcc[u^s,v^s]=\bigcap_{\al<\la}\ilcc[u_\al^s,v_\al^s]\neq\es.
\]
Since $f^{-1}(s)\neq\es$ for all $s\in\lexab$, $f$ is surjective.
Since $f$ is a monotone surjection and $f^{-1}(s)$ is a closed interval for all $s\in\lexab$, Proposition \ref{p:blex:2-1} implies that $f$ is continuous.

Let us show that $f$ is a homeomorphism; for this it suffices to check that $f$ is injective. Suppose otherwise. Then $|f^{-1}(s)|>1$, and hence $u^s<v^s$ for some $s\in\lexab$.
%
Denote $\te^s=\te_{u^s,v^s}$ and $\te_\al=\te_{u_\al^s,v_\al^s}$ for $\al<\la$.
Then
\[
\ilcc[u^s,v^s] \subsetneq \ilcc[u^s_\be,v^s_\be] \subsetneq \ilcc[u^s_\al,v^s_\al]
\qquad\text{ and } \qquad \te_\al\leq \te_\be\leq  \te^s < \la
\]
for $\al< \be < \la$. 
Hence either $u^s_\al<u^s$ for all $\al<\la$, or $v^s<v^s_\al$ for all $\al<\la$. Assume, without loss of generality, that $u^s_\al<u^s$ for all $\al<\la$.
Since $\la$ is a regular cardinal and the sequence $\sq{\te_\al}{\al<\la}$ is bounded and increasing, the sequence $\sq{\te_\al}{\al<\la}$ stabilizes: for some $\te<\la$ and $\al_0<\la$, $\te_\al=\te$ for all $\al$ with $\al_0\leq \al < \la$. Set
\[
C=\set{ \al<\la: \al_0<\al,\ u^s_{\al_0} < u^s_{\al} < u^s_{\al+1} }
\]
Then $|C|=\la$, $u^s_{\al} < a^s_\al< b^s_\al=u^s_{\al+1}$, and $(a^s_\al,b^s_\al)\in \cM_\te$ for $\al\in C$.
Set $A=\set{a^s_\al:\al\in C}$, $B=\set{b^s_\al:\al\in C}$, and $S=M_\te\cap(A\cup B)$. Since $\sset{a^s_\al,b^s_\al}\cap M_\te\neq\es$ for $\al\in C$, we get $u^s\in \cl S$ and $\chi(u^s,\cl S)=\la$. Hence $u^s\in \cl{M_\te}$ and $\chi(u^s,\cl{M_\te})=\la$, contradicting $\chi(\cl{M_\te})<\la$.
\end{proof}

\end{document}
